<+	+>	!comp!	!exe!
%        File: !comp!expand("%:p:t")!comp!
%     Created: !comp!strftime("%a %b %d %I:00 %p %Y ").substitute(strftime('%Z'), '\<\(\w\)\(\w*\)\>\(\W\|$\)', '\1', 'g')!comp!
% Last Change: !comp!strftime("%a %b %d %I:00 %p %Y ").substitute(strftime('%Z'), '\<\(\w\)\(\w*\)\>\(\W\|$\)', '\1', 'g')!comp!
%
\documentclass[UTF8, xcolor=table]{beamer}
%\usepackage{fontspec}
%\setsansfont{宋体}
\usepackage[BoldFont,SlantFont]{xeCJK}
%\setCJKmainfont[BoldFont={Adobe Heiti Std},ItalicFont={Adobe Kaiti Std}]{SimSun}

\usepackage{latexsym,amssymb,amsmath,amsbsy,amsopn,amstext,xcolor,multicol}
\usepackage{graphicx,wrapfig,fancybox}
\usepackage{pgf,pgfarrows,pgfnodes,pgfautomata,pgfheaps,pgfshade}
\usepackage{$HOME/.vim/ftplugin/latex-suite/packages/thubeamer}
%\usepackage[backend=bibtex,style=IEEE,sorting=none]{biblatex} % [参考文献格式](https://www.sharelatex.com/blog/2013/07/31/getting-started-with-biblatex.html)
\usepackage[backend=bibtex,sorting=none]{biblatex} % [参考文献格式](https://www.sharelatex.com/blog/2013/07/31/getting-started-with-biblatex.html) %mac IEEE not found
\usepackage{array}
\usepackage{bm}
\usepackage{caption}
\usepackage[caption=false,font=scriptsize]{subfig}
\usepackage{multirow}
\usepackage{booktabs}
\usepackage{tikz}
\usepackage{tikzscale}
\usepackage{animate}
%\usepackage{times} %与上面的冲突，加上这个 粗体斜体就失效
%\usepackage{mathptmx}


\defbibheading{bibliography}[\bibname]{} %avoid printbibliography 自动生成目录
\addbibresource{ref/papers-bib-in-google.bib}
\addbibresource{ref/chinese-ref.bib}
%\setbeamertemplate{bibliography item}{\insertbiblabel} %将列表中默认的丑陋的icon 改成数字,或者下面这个也行
\setbeamertemplate{bibliography item}[text] % [ref](http://tex.stackexchange.com/questions/68080/beamer-bibliography-icon)
%\setbeamertemplate{footline}[frame number]{}

%\setframeofframes{of}

\usepackage{boxedminipage} %for: bvh border
\def\fourgraphicswidth{0.35} %0.3\textwidth

\usepackage{algorithm} %%format of the algorithm
\usepackage{algpseudocode}
\floatname{algorithm}{算法}
\renewcommand{\algorithmicrequire}{\textbf{输入:}} %%Use Input in the format of Algorithm
\renewcommand{\algorithmicensure}{\textbf{输出:}} %%UseOutput in the format of Algorithm
%\algrenewcommand{\algorithmiccomment}[1]{\hskip3em $\rightarrow$ #1}
\algrenewcommand{\algorithmiccomment}[1]{$//$ #1}

\usepackage{listings}
\renewcommand\lstlistingname{代码}
\renewcommand\lstlistlistingname{代码}

\lstset{framexleftmargin=1.4em, xleftmargin=1.8em, basicstyle=\ttfamily\small, frame=single, numberstyle=\tiny, breaklines=true, keywordstyle=\color{blue!70}\bfseries,rulesepcolor=\color{red!20!green!20!blue!20}, numbers=none,fontadjust=true}
%frame=shadowbox, numberstyle=\tiny, breaklines=true,
%commentstyle=\color{red!50!green!50!blue!50},

\lstdefinelanguage{shader}{morekeywords={uniform, layout, uniform, vec2, vec3, vec4, in, out, gl_Position, dot, flat, int, float, gl_VertexID, xyz, w, x, y, z, location, version, sampler2DRect, bgr, gl_FragData, texture2DRect, gl_TexCoord, for, xy}, morecomment=[l]{//}}
%\setbeameroption{show notes} %un-comment to see the notes

%\usepackage{pgfpages}
%\renewcommand\pgfsetupphysicalpagesizes{%
%    \pdfpagewidth\pgfphysicalwidth\pdfpageheight\pgfphysicalheight%
%}
%\setbeameroption{show notes on second screen}

\begin{document}

\setbeamerfont{footnote}{size=\tiny}
\setbeamerfont{caption}{size=\scriptsize}
\setbeamertemplate{caption}[numbered]
\setbeamerfont{subsection in toc}{size=\footnotesize}
\renewcommand*{\bibfont}{\footnotesize}

\title{例会报告}
\author[陈达贵]{\vskip 20pt~~~~~~~学~~~~~~生：陈~~~~~达贵~~~~~~~~\vskip 5pt指导教师：陈~峰~教授}

\institute[清华大学~自动化系~\& 控制理论与技术研究所]{\small \vskip 38pt控制理论与技术研究所}
%\date{2015-06-07}
%\date{\small \vskip -17pt二〇一五年六月}
\date{\today}



\begin{frame}
\vspace{-15mm}
\titlepage{}
\vspace{-50mm}
\begin{figure}[htbp]
\begin{center}
\includegraphics[width=0.16\linewidth]{$HOME/.vim/ftplugin/latex-suite/packages/TsinghuaLogo.eps}
\end{center}
\end{figure}
%\beign{picture}(1,1)
%\put(6,8){\includegraphics[width=0.15\linewidth]{Tsinghua_University_Logo.eps}}
%\end{picture}
\end{frame}

\section{Outline}
\frame{
\frametitle{Outline}
\begin{enumerate}
\item Doom相关工作：
\item 科研工作:
\item IT相关工作：
\item 专利相关工作：
\end{enumerate}
}

\section{Doom工作}
\frame{
\frametitle{<++>}
\begin{enumerate}
\item <++>
\end{enumerate}
}

\section{科研工作}
\begin{frame}[shrink]
\frametitle{<++>}
\begin{enumerate}
\item <++>
\end{enumerate}
\end{frame}

\section{下周计划}

\begin{frame}
\frametitle{下周计划}
<++>
\end{frame}

\section{文献分享}
\begin{frame}{文献分享}
\includegraphics[width=\linewidth]{<++>}
\end{frame}

\end{document}
